\documentclass[11pt]{article}

\usepackage[a4paper,margin=30mm]{geometry}
\usepackage[T1]{fontenc}
\usepackage[utf8]{inputenc}
\usepackage{lmodern}
\usepackage{microtype}
\usepackage{amsmath,amssymb,amsthm,mathtools}
\usepackage{booktabs}
\usepackage{graphicx}
\usepackage{xcolor}
\usepackage[hidelinks]{hyperref}
\usepackage[nameinlink,noabbrev]{cleveref}

\hypersetup{
  pdftitle={Window-covering words and a bouquet of cycles},
  pdfauthor={Nicolo Pandolfo},
  pdfsubject={Exact enumeration of OEIS A242239},
  pdfkeywords={combinatorics on words, finite automata, generating functions, OEIS A242239}
}

\newtheorem{theorem}{Theorem}[section]
\newtheorem{proposition}[theorem]{Proposition}
\newtheorem{lemma}[theorem]{Lemma}
\newtheorem{corollary}[theorem]{Corollary}
\theoremstyle{remark}
\newtheorem{remark}[theorem]{Remark}

\newcommand{\SigmaK}{\Sigma_k}
\newcommand{\W}{\mathcal W}
\newcommand{\B}{\mathcal B}
\newcommand{\Rstate}{\mathsf R}
\newcommand{\GF}{\mathcal A}

\title{\textbf{Window-covering words and a bouquet of cycles}\\
\large Exact enumeration of the OEIS family A242239}
\author{Nicol\`o Pandolfo}
\date{July 23, 2026\\\small Draft for independent mathematical review}

\begin{document}
\maketitle

\begin{abstract}
For integers $k,n\geq 1$, let $T(n,k)$ be the number of canonical words of length
$n+k+1$ over $\{0,1,\ldots,k\}$ such that every consecutive block of length
$k+2$ contains every symbol.  This is the table OEIS A242239.  We prove that
these words, modulo relabeling, are encoded by length-$n$ walks on a directed
bouquet whose cycle lengths are exactly $1,2,\ldots,k+1$.  Consequently,
\[
 \sum_{n\geq 1}T(n,k)x^n
 =\frac{\displaystyle\sum_{j=1}^{k+1}
       \frac{(k+2-j)(kj+2)}{2}x^j}
      {1-x-x^2-\cdots-x^{k+1}}.
\]
This proves the empirical $(k+1)$-bonacci recurrence recorded for every column
of A242239, with exact starting index $n=k+2$, and also proves the stable-row
formula
\[
 T(n,k)=\frac{2^n-1}{2}k^2+\frac{2n+1}{2}k+1
 \qquad (k\geq n-3).
\]
For $k=4$ the result proves the empirical recurrence and generating function
of A242235.  Using the already known generating function of A145113, it also
proves Mathar's conjectured identity for that sequence.
\end{abstract}

\section{Introduction and main result}

Set $\SigmaK=\{0,1,\ldots,k\}$, so $|\SigmaK|=k+1$.  A word over
$\SigmaK$ is called \emph{canonical} if the first occurrences of its distinct
symbols appear in the order $0,1,2,\ldots$; equivalently, it is a restricted
growth word.  The table A242239 is defined by
\begin{equation}\label{eq:defT}
 T(n,k)=\#\left\{
 w\in\SigmaK^{n+k+1}:
 \begin{array}{l}
 w\text{ is canonical, and}\\
 \text{every }(k+2)\text{-factor of }w\text{ contains }\SigmaK
 \end{array}
 \right\}.
\end{equation}
The OEIS entry records, as empirical observations, both a $(k+1)$-step
Fibonacci recurrence for each column and a quadratic formula for each row once
$k$ is sufficiently large \cite{oeisA242239}.  The special column $k=4$ is
A242235; its recurrence and generating function are likewise marked empirical,
and an identity proposed by Mathar is marked conjectural \cite{oeisA242235}.

The main result gives the complete ordinary generating function in closed form.

\begin{theorem}[Exact generating function]\label{thm:main}
For every $k\geq 1$, define
\[
 D_k(x)=1-x-x^2-\cdots-x^{k+1}
\]
and
\[
 N_k(x)=\sum_{j=1}^{k+1}c_{k,j}x^j,
 \qquad
 c_{k,j}=\frac{(k+2-j)(kj+2)}{2}.
\]
Then
\begin{equation}\label{eq:mainGF}
 \GF_k(x):=\sum_{n\geq1}T(n,k)x^n=\frac{N_k(x)}{D_k(x)}.
\end{equation}
\end{theorem}

The coefficients $c_{k,j}$ are positive integers.  The integrality is immediate:
if $k$ or $j$ is even then $kj+2$ is even, while if both are odd then
$k+2-j$ is even.

\begin{corollary}[Column recurrence]\label[corollary]{cor:recurrence}
For every $k\geq1$,
\[
 T(n,k)=\sum_{i=1}^{k+1}T(n-i,k)
 \qquad (n\geq k+2).
\]
This range is exact: with the convention $T(m,k)=0$ for $m\leq0$, the same
relation fails for every $1\leq n\leq k+1$, because its residual is the positive
coefficient $c_{k,n}$.
\end{corollary}

\begin{corollary}[Stable rows]\label[corollary]{cor:rows}
For $n,k\geq1$ with $k\geq n-3$,
\begin{equation}\label{eq:rowformula}
 T(n,k)=\frac{2^n-1}{2}k^2+\frac{2n+1}{2}k+1.
\end{equation}
\end{corollary}

As of the access date in the bibliography, these formulas remain labelled
empirical on A242239 and its relevant column entries.  We are not aware of a
previous proof; this statement is a qualified literature-status report, not a
claim that an exhaustive priority search is possible.

\section{From words to a defect-pair automaton}

Put $q=k+1$.  A block of length $q$ is \emph{rainbow} if all its letters are
distinct.  Otherwise, a valid block will have one repeated symbol.  If its two
copies occur in positions $1\leq i<j\leq q$, call its type $(i,j)$.

\begin{lemma}[Relabeling orbits]\label[lemma]{lem:orbits}
Let $N\geq q+1$, and let $\W_k(N)$ be the set of labelled words of length $N$
over $\SigmaK$ satisfying the sliding-window condition in \eqref{eq:defT}.
The symmetric group $S_q$ acts freely on $\W_k(N)$ by relabeling, and every
orbit contains exactly one canonical word.  Hence
\[
 T(n,k)=\frac{|\W_k(n+q)|}{q!}.
\]
\end{lemma}

\begin{proof}
The first window of length $q+1$ contains every symbol, so every word uses all
$q$ symbols.  A permutation fixing the word therefore fixes every symbol and
is the identity.  In each orbit, relabel the symbols by their first-occurrence
order; this produces exactly one canonical representative.
\end{proof}

\begin{lemma}[Block classification]\label[lemma]{lem:block}
Every length-$q$ block occurring in a word of $\W_k(N)$ contains at least
$q-1$ distinct symbols.  It is therefore either rainbow or has a unique type
$(i,j)$ and misses exactly one symbol.
\end{lemma}

\begin{proof}
Any length-$q$ block extends to a valid length-$(q+1)$ window, to the right
unless it is the final block and to the left otherwise.  Deleting one letter
from a window containing all $q$ symbols leaves at least $q-1$ distinct symbols.
A non-rainbow block consequently contains exactly one doubled symbol and misses
exactly one symbol.
\end{proof}

\begin{lemma}[Transitions]\label[lemma]{lem:transitions}
Let $B=b_1\cdots b_q$ be the final block of a valid prefix.
\begin{enumerate}
\item[\textup{(a)}] If $B$ is rainbow, every symbol may be appended.  Appending $b_1$
      gives another rainbow block, while appending $b_{p+1}$ gives type
      $(p,q)$ for $1\leq p\leq q-1$.
\item[\textup{(b)}] If $B$ has type $(i,j)$ and misses $m$, the only legal appended symbol
      is $m$.  The new block is rainbow when $i=1$, and has type
      $(i-1,j-1)$ when $i>1$.
\end{enumerate}
Each listed transition is realized by exactly one appended symbol.
\end{lemma}

\begin{proof}
The newly created window is $Bx$, and it is valid precisely when the symbols
in $B$ together with $x$ cover $\SigmaK$.  In the rainbow case, shifting removes
$b_1$ and appending a copy of $b_r$ either replaces the removed symbol ($r=1$)
or creates a doubled symbol in positions $(r-1,q)$.  In the defect case, the
missing symbol is forced; shifting either removes the first copy of the doubled
symbol or shifts both copies one position to the left.
\end{proof}

The quotient graph has a rainbow state $\Rstate$ and one state $(i,j)$ for each
$1\leq i<j\leq q$.  The transitions in \cref{lem:transitions} reveal a simpler
structure.  For each $\ell=2,\ldots,q$, the states
\[
 (\ell-1,q),(\ell-2,q-1),\ldots,(1,q-\ell+2)
\]
form, together with $\Rstate$, a directed cycle of length $\ell$.  Different
values of $\ell$ give disjoint non-hub vertices, and $\Rstate$ also has a
self-loop.  Thus the quotient graph is the directed bouquet
\[
 \B_q=C_1\vee C_2\vee\cdots\vee C_q,
\]
where the cycles share only the hub $\Rstate$.

\begin{proposition}[Bouquet encoding]\label[proposition]{prop:bouquet}
For every $n,k\geq1$, $T(n,k)$ is the number of directed walks of length $n$ in
$\B_{k+1}$, with every vertex allowed as a starting vertex and every vertex
allowed as an ending vertex.
\end{proposition}

\begin{proof}
Choose the initial block of length $q$ and record its type.  By
\cref{lem:transitions}, a state walk determines each subsequently appended
letter uniquely, and every state edge is realizable.  Hence labelled words are
in bijection with pairs consisting of a concrete initial block and a state walk
of length $n$.

For each state type, the set of concrete blocks has exactly one $S_q$-orbit.
There are $q!$ rainbow blocks.  For a fixed defect type $(i,j)$, choose the
doubled symbol and then injectively assign the remaining $q-2$ positions from
the other $q-1$ symbols; again there are $q!$ blocks.  Relabeling changes the
concrete block but not the state walk.  Applying \cref{lem:orbits} therefore
leaves one initial orbit for every vertex of $\B_q$, proving the result.
\end{proof}

\section{Enumeration of walks on the bouquet}

Let
\begin{equation}\label{eq:H}
 H_k(x)=1+\sum_{j=1}^{k}(k+1-j)x^j,
 \qquad
 s_k=1+\frac{k(k+1)}2=\frac{k^2+k+2}{2}.
\end{equation}
The number $s_k$ is the number of vertices of $\B_{k+1}$.

Let $P_k(x)$ be the generating function for walks that start at the hub
$\Rstate$ and may end anywhere.  Such a walk consists uniquely of a sequence of
complete returns to the hub, followed by a proper prefix of one cycle.  The
complete return lengths are $1,2,\ldots,k+1$, one of each length.  A proper
prefix of length $j\geq1$ can be chosen in $k+1-j$ ways.  Therefore
\begin{equation}\label{eq:Pk}
 P_k(x)=\frac{H_k(x)}{D_k(x)}.
\end{equation}

A vertex at forced distance $i$ from the hub contributes
\[
 1+x+\cdots+x^{i-1}+x^iP_k(x),
\]
and there are $k+1-i$ such vertices.  Summing over all vertices, including the
hub, gives the generating function for walks of all lengths:
\begin{equation}\label{eq:walkall}
 \sum_{n\geq0}T_0(n,k)x^n
 =H_k(x)P_k(x)+G_k(x),
\end{equation}
where $T_0(0,k)=s_k$, $T_0(n,k)=T(n,k)$ for $n\geq1$, and
\[
 G_k(x)=\sum_{i=1}^{k}(k+1-i)(1+x+\cdots+x^{i-1})
       =\frac{s_k-H_k(x)}{1-x}.
\]
The elementary identity
\begin{equation}\label{eq:keyidentity}
 (1-x)H_k(x)=D_k(x)+kx
\end{equation}
now removes the only apparent extra denominator.  From
\eqref{eq:Pk}--\eqref{eq:walkall},
\begin{align*}
 \GF_k(x)
 &=\frac{H_k(x)^2}{D_k(x)}+\frac{s_k-H_k(x)}{1-x}-s_k\\
 &=\frac{H_k(x)^2}{D_k(x)}+\frac{s_kx-H_k(x)}{1-x}.
\end{align*}
Consequently,
\begin{align}
 D_k(x)\GF_k(x)
 &=\frac{(1-x)H_k(x)^2+D_k(x)(s_kx-H_k(x))}{1-x}\notag\\
 &=\frac{x\bigl(kH_k(x)+s_kD_k(x)\bigr)}{1-x}.
 \label{eq:Nintermediate}
\end{align}
The polynomial inside the last quotient vanishes at $x=1$, since
$H_k(1)=s_k$ and $D_k(1)=-k$.  Thus the division by $1-x$ is exact.

For completeness, let $N_k(x)=\sum_{j=1}^{k+1}c_{k,j}x^j$.  Dividing
\eqref{eq:Nintermediate} by taking partial sums of coefficients gives
\begin{align*}
 c_{k,j}
 &=k+s_k+\sum_{r=1}^{j-1}\bigl(k(k+1-r)-s_k\bigr)\\
 &=\frac{(k+2-j)(kj+2)}2,
 \qquad 1\leq j\leq k+1.
\end{align*}
This proves \cref{thm:main}, and \cref{cor:recurrence} follows by coefficient
comparison in $D_k(x)\GF_k(x)=N_k(x)$.

\begin{table}[ht]
\centering
\small
\caption{The first numerator polynomials, generated directly from
$c_{k,j}=((k+2-j)(kj+2))/2$.}
\label{tab:numerators}
\resizebox{\textwidth}{!}{\begin{tabular}{c l}
\toprule
$k$ & $N_k(x)$ \\
\midrule
1 & $3x + 2x^{2}$ \\
2 & $6x + 6x^{2} + 4x^{3}$ \\
3 & $10x + 12x^{2} + 11x^{3} + 7x^{4}$ \\
4 & $15x + 20x^{2} + 21x^{3} + 18x^{4} + 11x^{5}$ \\
5 & $21x + 30x^{2} + 34x^{3} + 33x^{4} + 27x^{5} + 16x^{6}$ \\
6 & $28x + 42x^{2} + 50x^{3} + 52x^{4} + 48x^{5} + 38x^{6} + 22x^{7}$ \\
7 & $36x + 56x^{2} + 69x^{3} + 75x^{4} + 74x^{5} + 66x^{6} + 51x^{7} + 29x^{8}$ \\
8 & $45x + 72x^{2} + 91x^{3} + 102x^{4} + 105x^{5} + 100x^{6} + 87x^{7} + 66x^{8} + 37x^{9}$ \\
\bottomrule
\end{tabular}
}
\end{table}

\section{The stable-row formula}

We prove \cref{cor:rows} directly from the bouquet.  Let
\[
 r_m=[x^m]P_k(x)
\]
be the number of length-$m$ walks starting at the hub.  From
$D_k(x)P_k(x)=H_k(x)$, a short induction gives
\begin{equation}\label{eq:hubcounts}
 r_0=1,
 \qquad
 r_m=k2^{m-1}+1\quad(1\leq m\leq k+2),
\end{equation}
and at the next index
\begin{equation}\label{eq:hubboundary}
 r_{k+3}=k2^{k+2}-k+1.
\end{equation}
Indeed, for $1\leq m\leq k$, the recurrence for $r_m$ has forcing term
$k+1-m$; summing the earlier values in \eqref{eq:hubcounts} gives the claimed
value.  The cases $m=k+1,k+2,k+3$ then follow from the homogeneous recurrence
with the appropriate $k+1$ preceding terms.

A vertex at distance $i$ from the hub has exactly one length-$n$ continuation
if $n\leq i$, and otherwise reaches the hub after $i$ steps and has
$r_{n-i}$ continuations.  Therefore
\begin{equation}\label{eq:rowwalksum}
 T(n,k)=r_n+\sum_{i=1}^{k}(k+1-i)
 \begin{cases}
  1,& n\leq i,\\
  r_{n-i},& n>i.
 \end{cases}
\end{equation}
If $n\leq k$, substituting \eqref{eq:hubcounts} into
\eqref{eq:rowwalksum} gives
\begin{align*}
 T(n,k)
 &=k2^{n-1}+1
   +\sum_{i=1}^{n-1}(k+1-i)(k2^{n-i-1}+1)
   +\sum_{i=n}^{k}(k+1-i)\\
 &=\frac{2^n-1}{2}k^2+\frac{2n+1}{2}k+1.
\end{align*}
For $n=k+1$ and $n=k+2$, the same substitution applies with the final sum
empty.  For $n=k+3$, use \eqref{eq:hubboundary} for the first term and
\eqref{eq:hubcounts} for every term in the remaining sum.  The same
simplification results.  These four cases are exactly $n\leq k+3$, proving
\cref{cor:rows}.

\section{The column A242235 and Mathar's identity}

Taking $k=4$ in \cref{thm:main} gives
\begin{equation}\label{eq:A242235GF}
 \sum_{n\geq1}T(n,4)x^n
 =\frac{x(15+20x+21x^2+18x^3+11x^4)}
 {1-x-x^2-x^3-x^4-x^5}.
\end{equation}
Thus A242235 satisfies the pentanacci recurrence from $n=6$, with initial
values
\[
 15,35,71,139,271.
\]
This proves the empirical formulas recorded on the OEIS entry
\cite{oeisA242235}.

Let $b(m)$ be A145113, the number of length-$m$ binary words with fewer than
five zeros between consecutive ones.  Its generating function is already
recorded on that entry as
\begin{equation}\label{eq:A145113GF}
 B(x)=\sum_{m\geq0}b(m)x^m
 =\frac{1-x+x^6}{1-3x+2x^2+x^6-x^7}
 =\frac{1-x+x^6}{(1-x)^2D_4(x)}.
\end{equation}
See \cite{oeisA145113}.  It can also be obtained from the standard seven-state
automaton that records the number of zeros since the most recent one.

Combining \eqref{eq:A242235GF} and \eqref{eq:A145113GF} and clearing the common
denominator gives the exact identity
\begin{equation}\label{eq:MatharGF}
 \GF_4(x)=16xB(x)+\frac{4x}{(1-x)^2}-\frac{5x}{1-x}.
\end{equation}
Hence
\begin{equation}\label{eq:Mathar}
 T(n,4)=16b(n-1)+4n-5\qquad(n\geq1).
\end{equation}
For $n\geq2$, this is precisely Mathar's conjecture
$T(n,4)=16\,\mathrm{A145113}(n-1)+\mathrm{A004767}(n-2)$, because
A004767$(m)=4m+3$ \cite{oeisA004767}.  At $n=1$, \eqref{eq:Mathar} remains
valid, but the A004767 notation requires polynomially extending $4m+3$ to
$m=-1$; it is not an ordinary in-range OEIS sequence evaluation.

\section{Scope, verification, and disclosure}

The bouquet reduction depends critically on the window length being exactly
one larger than the alphabet size.  In that regime, deleting one letter from a
covering window leaves either a rainbow block or a block with one doubled and
one missing symbol.  For longer windows, multiple defects can occur, and the
same state space does not apply without further work.

No step of the proof relies on numerical evidence.  The accompanying verification program is
\path{verification/verify_hardin_family.py}.  It uses only the Python standard library and
independently checks:
\begin{enumerate}
\item the closed numerator formula against walks on the bouquet;
\item a raw last-block DFA constructed only from the defining window condition;
\item direct enumeration of canonical restricted-growth words for small cases;
\item the first $10\times8$ values displayed on A242239;
\item the column recurrence, stable-row formula, A242235, A145113, and the exact
      polynomial identity underlying \eqref{eq:MatharGF}.
\end{enumerate}
The script also generates \cref{tab:numerators}, preventing manual
transcription from becoming part of the verification chain.

AI systems were used during exploration, adversarial checking, and drafting.
The mathematical deliverable is the explicit proof above together with a
deterministic, inspectable verification program; no AI output is treated as an
oracle or as evidence by itself.

\section*{Acknowledgment}
The definitions, numerical data, and empirical conjectures originate with the
contributors to the cited OEIS entries, in particular R. H. Hardin, R. J.
Mathar, and Colin Barker.

\begin{thebibliography}{9}

\bibitem{oeisA242239}
OEIS Foundation Inc.,
\emph{A242239: Number of length $n+k+1$ arrays with every value appearing in
every consecutive $k+2$ elements},
\url{https://oeis.org/A242239}, accessed July 23, 2026.

\bibitem{oeisA242235}
OEIS Foundation Inc.,
\emph{A242235},
\url{https://oeis.org/A242235}, accessed July 23, 2026.

\bibitem{oeisA145113}
OEIS Foundation Inc.,
\emph{A145113: Binary words with fewer than five zeros between consecutive
ones},
\url{https://oeis.org/A145113}, accessed July 23, 2026.

\bibitem{oeisA004767}
OEIS Foundation Inc.,
\emph{A004767: $a(n)=4n+3$},
\url{https://oeis.org/A004767}, accessed July 23, 2026.

\bibitem{stanley}
R. P. Stanley,
\emph{Enumerative Combinatorics, Volume 1}, 2nd ed.,
Cambridge University Press, 2012, Section 4.7.

\bibitem{flajolet}
P. Flajolet and R. Sedgewick,
\emph{Analytic Combinatorics},
Cambridge University Press, 2009, Chapter I.

\end{thebibliography}

\end{document}
